Space optical communication research has concentrated overwhelmingly on
\emph{trunk} links: high-rate, long-range channels between spacecraft, or
between a spacecraft and a ground station. This section reviews that literature
and then the areas relevant to the opposite regime of short-range proximity
links between co-located surface assets: short-range infrared signalling,
multi-detector receiver front ends, maximum-likelihood sequence detection, and
coding for channels that drop symbols. The last two supply the method adopted
here; the first three establish the link model and the gap this work addresses.
A closing subsection covers optical localization, the platform's original
objective, which is motivated but not evaluated in this paper.

\subsection{Space Optical Communication}

Prior work on free-space optical links establishes the performance benchmarks
and design principles achieved by intensity-modulation direct-detection (IM/DD)
systems across a wide range of distances. Hemmati~\cite{Hemmati2006} provides a
systematic treatment of deep-space optical link architectures, deriving received
power margins and pointing loss budgets as functions of aperture, beam
divergence, and range. Biswas and Piazzolla~\cite{Biswas2003} quantified these
parameters for a Mars-to-Earth downlink, developing link-budget analyses based
on received optical power and noise processes. At orbital distances,
Sodnik~et~al.~\cite{Sodnik2010} demonstrated bidirectional optical
communication in the European Space Agency SILEX program, validating
acquisition, pointing, and tracking across an intersatellite link.
Boroson~et~al.~\cite{Boroson2014} reported the Lunar Laser Communication
Demonstration, achieving a 622~Mbps downlink from lunar orbit using a laser
transmitter and photon-counting receiver.

These systems share a common design posture. Range is large, the link budget is
the binding constraint, and the receiver is a substantial instrument in its own
right: cryogenically cooled detectors, precision gimbals, dedicated signal
processing hardware. Performance is won through aperture, transmit power, and
pointing accuracy.

Proximity links on a planetary surface invert every one of those assumptions.
Range is metres, so link budget is comfortable. The receiver is a photodiode
and a shared microcontroller on a mass- and power-constrained mobile asset.
What limits the link is not how much optical power arrives but how much of the
arriving signal a cheap receiver can successfully decode. The present work
addresses that regime, which the trunk-link literature does not speak to.

\subsection{Short-Range Infrared Links}

Short-range optical wireless communication using light-emitting diodes has been
characterized extensively for terrestrial indoor use, and that literature
supplies the channel and receiver models applicable to surface proximity links.
Gfeller and Bapst~\cite{Gfeller1979} established diffuse infrared as a
practical in-building medium. Kahn and Barry~\cite{Kahn1997} provided the
foundational analysis of indoor infrared channels, deriving the tradeoffs among
transmit power, receiver sensitivity, ambient-light shot noise, thermal noise,
and achievable performance. Their treatment of IM/DD signalling, including
pulse-position and pulse-width formats in which information is carried in the
time domain rather than in amplitude, directly motivates the modulation used
here. Ghassemlooy~et~al.~\cite{Ghassemlooy2012} present a comprehensive
treatment of channel modelling, path loss, multipath effects, and photodetector
noise. Elgala~et~al.~\cite{Elgala2011} survey the resulting system-level
tradeoffs, including the tension among receiver field-of-view, angular
sensitivity, and ambient-light rejection that governs photodiode front-end
design.

Ambient-light rejection dominates much of this terrestrial work, and it is
precisely the constraint that relaxes in a permanently shadowed lunar region,
where solar infrared background is largely absent. What does \emph{not} relax
is the receiver's own limitation. Time-domain symbol encoding is the relevant
inheritance: where an analog front end operates outside its linear region (as
is measured to be the case for the platform described here), amplitude carries
no usable information while symbol timing remains intact. The same argument
applies to a surface deployment for different reasons, since regolith dust
accumulation and thermal cycling degrade optical throughput and detector
responsivity over a mission without disturbing symbol timing.

\subsection{Multi-Detector Receivers}

Receivers employing more than one photodetector have been studied primarily for
multipath and ambient-light rejection. Carruthers and Kahn~\cite{Carruthers2000}
analyse angle-diversity receivers in which multiple detectors with distinct
orientations are combined, showing improved performance in non-directed links
where a single wide field-of-view detector suffers co-channel interference and
multipath dispersion. The receiver described here uses two photodiodes for a
different purpose: detection robustness at the pulse level, with edge detection
driven by the summed response so that a symbol arriving predominantly on one
detector is not lost. On a mobile surface asset, where relative pointing between
rover and lander changes continuously and cannot be assumed stable, this
tolerance to partial illumination is a practical requirement rather than an
optimisation. The normalized difference between channels, which cancels terms
common to both, is noted as a bearing observable but is not evaluated here.

\subsection{Maximum-Likelihood Sequence Detection}

Recovering a transmitted sequence by searching for the most likely path through
a trellis, rather than deciding each symbol independently, was formalized by
Forney~\cite{Forney1972} for channels with intersymbol interference. The
underlying dynamic program is due to Viterbi~\cite{Viterbi1967}, introduced for
decoding convolutional codes and later given a general treatment by
Forney~\cite{Forney1973}. The same recursion decodes hidden Markov models,
whose canonical exposition is Rabiner's~\cite{Rabiner1989}: the observations are
noisy and the underlying state sequence is unobserved.

A standard hidden Markov model pairs each state transition with exactly one
observation, and it is that correspondence which the link measured here
violates. Two established families relax it. Explicit-duration, or hidden
semi-Markov, models allow a state to persist for a sampled duration rather than
a single step~\cite{Rabiner1989}. Profile hidden Markov models, developed for
biological sequence alignment, introduce explicit \emph{delete} states that
advance the chain while consuming no observation at all~\cite{Krogh1994}. The
decoder described in Section~\ref{system_overview} is structurally of the
latter kind: a pulse the receiver never detects still advances the byte
framing, but produces nothing for the model to condition on.

That property is what makes the approach applicable here. The failure mode
measured on this link is not symbol confusion but symbol \emph{omission}: the
receiver fails to detect a pulse, so no observation is produced for it. A
per-symbol classifier has no mechanism to recover such a symbol, whereas a
sequence model can represent it as an unobserved state transition and infer it
from timing evidence and structural constraints.

\subsection{Deletion Channels and Synchronization Errors}

A channel that drops symbols without announcing it is, in coding-theoretic
terms, a \emph{deletion} channel, and codes for channels with insertion and
deletion errors have been studied since Levenshtein~\cite{Levenshtein1966};
Mercier et al.~\cite{Mercier2010} survey the field. The closest prior work to
the decoder used here is Davey and MacKay's watermark
code~\cite{Davey2001}, which embeds a known pseudorandom sequence in the
transmitted stream and recovers synchronization with an inner hidden Markov
decoder whose state space contains explicit insertion and deletion transitions.
The mechanism this paper uses is therefore not new. A hidden Markov model with
deletion transitions, decoded by dynamic programming, is the established tool
for this channel, and the present decoder should be understood as an instance
of it rather than as a new construction.

Two things differ, and both follow from the physical layer rather than from the
algorithm. First, that literature generally assumes the transmitter can be
designed: watermark and marker codes work by inserting redundancy for the
receiver to synchronize against. Here the transmitter is fixed and nothing is
added. The structure the decoder exploits, byte framing and a restricted
payload alphabet, is pre-existing protocol overhead rather than deliberately
inserted synchronization information, which is what makes the approach
available as a retrofit to a deployed or mass-constrained system where the
transmitter cannot be changed.

Second, and more consequentially, a deletion in the classical model is
invisible. The receiver observes a shorter sequence with no indication of where
the loss occurred, and supplying that missing information is precisely why a
watermark is inserted. In a pulse-width format the deletion is not invisible:
it lengthens the following inter-pulse interval by a known, symbol-dependent
amount, so the analog timing of the physical layer supplies for free the
synchronization evidence a watermark exists to provide. A symbol-level deletion
channel model discards exactly the quantity that makes recovery possible here.

The contribution of this work is accordingly empirical rather than theoretical.
It reports a measurement, on hardware, that the timing side-information carried
by a pulse-width physical layer is present in every single-omission
transmission observed and is sufficient to localize and identify the lost
symbol without added redundancy, together with a quantification of how much of
the resulting message loss a sequence decoder actually recovers. Sequence
detection is already well established in space communication in the form of
convolutional coding with Viterbi decoding; what is measured here is what the
same machinery buys when no redundancy is available to decode.

\subsection{Optical Localization}

Received optical signal strength has been investigated as a position observable
where emitter positions are known \textit{a priori}. Do and Yoo~\cite{DoYoo2016}
survey visible-light positioning based on signal strength, angle of arrival, and
time difference of arrival, and Lim~\cite{Lim2015} demonstrates
three-dimensional LED-based positioning with sub-decimeter accuracy under
line-of-sight conditions. Such techniques are attractive for surface operations
because they would reuse an existing communication link as a navigation aid,
avoiding dedicated sensing hardware in environments where GNSS is unavailable
and visual odometry is degraded by low-contrast regolith and harsh shadowing.
They are not evaluated in this paper: measurements reported in
Section~\ref{results} show the receiver's amplitude response to be flat with
distance over the range tested, which precludes signal-strength ranging on the
present front end and is addressed as future work.
